\begin{mainquestion}{capacitortimeconstant}{Norton Hears a Who!}

\questionpoints
The circuit below is driven by a DC current source $I_S$. The switch $S_1$ has been \textbf{open} for a long time. At time $t=0$, the switch $S_1$ is \textbf{closed}.

\begin{center}
\begin{circuitikz}[american, scale=1.2]
    % Source Branch
    \draw (0,0) to[I, l=$I_S{=}10\text{mA}$] (0,3);
    
    % Bottom Rail
    \draw (0,0) -- (8,0);
    \draw (4,0) node[ground]{}; 
    
    % Top Rail
    \draw (0,3) -- (8,3);
    \draw (4,3) node[above] {Node A} node[circ]{};

    % R1 Branch (Always connected)
    \draw (2,3) to[R, l=$R_1{=}4\text{k}\Omega$] (2,0);
    
    % Switched Branch R2
    % (nos = normally open switch)
    \draw (4,3) to[nos, l=$S_1$] (4,1.5)
                to[R, l=$R_2{=}4\text{k}\Omega$] (4,0);

    % Capacitor Branch (Right side)
    % R3 and C in series
    \draw (8,3) to[R, l=$R_3{=}2\text{k}\Omega$] (8,1.5)
                to[C, l=$C{=}5\mu\text{F}$, v=$v_c(t)$] (8,0);
\end{circuitikz}
\end{center}

% =========================================================
\begin{partquestion}{initialvoltage}
{Calculate the initial voltage across the capacitor, $v_c(0^-)$.}

\begin{answer}[2in]
        \textbf{State $t < 0$ (Switch Open):}
        Since the switch is open, $R_2$ is disconnected. In DC steady state, the capacitor acts as an open circuit. Therefore, all current $I_S$ flows through $R_1$.
        
        The voltage at Node A is:
        \[ V_A = I_S \cdot R_1 = 10\text{mA} \cdot 4\text{k}\Omega = 40\text{V} \]
        
        Since no current flows through $R_3$ (open circuit branch), there is no voltage drop across it.
        \[ v_c(0^-) = V_A = 40\text{V} \]
\end{answer}

\end{partquestion}

% =========================================================
\begin{partquestion}{switchcurrent}
{Calculate the current flowing through the switch $S_1$ immediately after it closes, at $t=0^+$.}

\begin{answer}[2in]
        At $t=0^+$, the capacitor acts as a voltage source with value $v_c(0^-) = 40\text{V}$.
        Let $V_A(0^+)$ be the voltage at Node A at this instant.
        Apply KCL at Node A:
        \[ I_S = i_{R1} + i_{sw} + i_{cap} \]
        \[ I_S = \frac{V_A(0^+)}{R_1} + \frac{V_A(0^+)}{R_2} + \frac{V_A(0^+) - v_c(0^+)}{R_3} \]
        
        Substitute values ($I_S=10\text{mA}, R_1=4\text{k}, R_2=4\text{k}, R_3=2\text{k}, v_c=40\text{V}$):
        \[ 10 = \frac{V_A}{4} + \frac{V_A}{4} + \frac{V_A - 40}{2} \]
        Multiply by 4 to clear denominators:
        \[ 40 = V_A + V_A + 2(V_A - 40) \]
        \[ 40 = 2V_A + 2V_A - 80 \implies 4V_A = 120 \implies V_A(0^+) = 30\text{V} \]
        
        The current through the switch is:
        \[ i_{sw}(0^+) = \frac{V_A(0^+)}{R_2} = \frac{30\text{V}}{4\text{k}\Omega} = 7.5\text{mA} \]
\end{answer}

\qnewpage
\end{partquestion}

% =========================================================
\begin{partquestion}{norton}
{Treating the capacitor as the load, calculate the \textbf{Norton Equivalent Circuit} ($I_{N}$, $R_{N}$) of the remaining network for $t > 0$.}

\begin{answer}[3in]
        \textbf{State $t > 0$ (Switch Closed):}
        $R_2$ is now in parallel with $R_1$.
        
        \textbf{1. Norton Resistance ($R_N$):}
        Open the current source ($I_S = 0$). Looking back from the capacitor terminals, we see $R_3$ in series with ($R_1 || R_2$).
        \[ R_{12} = 4\text{k} || 4\text{k} = 2\text{k}\Omega \]
        \[ R_N = R_3 + R_{12} = 2\text{k} + 2\text{k} = 4\text{k}\Omega \]

        \textbf{2. Norton Short Circuit Current ($I_N$):}
        Short the capacitor terminals to ground. $I_S$ splits between $R_{12}$ (the parallel combination of $R_1, R_2$) and the path containing $R_3$.
        Since $R_{12} = 2\text{k}\Omega$ and $R_3 = 2\text{k}\Omega$, the current splits equally.
        \[ I_N = I_S \cdot \frac{1}{2} = 5\text{mA} \]
\end{answer}

\end{partquestion}

% =========================================================
\begin{partquestion}{voltageexpression}
{Using the Norton equivalent circuit you found in part (c), write a mathematical expression for the voltage across the capacitor $v_c(t)$ for $t \ge 0$.}

\begin{answer}[2in]
        We use the step response equation: $v_c(t) = v(\infty) + [v(0) - v(\infty)]e^{-t/\tau}$.
        
        \textbf{Final Value:}
        Using the Norton equivalent:
        \[ v_c(\infty) = I_N \cdot R_N = 5\text{mA} \cdot 4\text{k}\Omega = 20\text{V} \]
        
        \textbf{Time Constant:}
        \[ \tau = R_N \cdot C = 4\text{k}\Omega \cdot 5\mu\text{F} = 20\text{ms} = 0.02\text{s} \]
        
        \textbf{Expression:}
        \[ v_c(t) = 20 + (40 - 20)e^{-t/0.02} \]
        \[ v_c(t) = 20 + 20e^{-50t} \text{ Volts} \]
\end{answer}

\end{partquestion}

\end{mainquestion}
